% page 596 du fac-simile (image p-595 du scan)
% bloc 157.5 x 246.3 mm ; marge gauche 8.0 mm
\pgc{-12.700mm}{0.000mm}{
\l{\cel{3}588.}
\l{}
\l{\sou{tragakanto}. (\sou{bot.})\cel{2}Arboreto ek la genero \textquotedbl{}astragalo\textquotedbl{}, qua}
\l{\cel{3}donas adraganto. - L.\cel{1}\sou{astragalus}. DEFR.}
\l{}
\l{\sou{tragedio}. I. Teatrajo normala, di qua la roli luktas kontre}
\l{\cel{3}sua pasioni o kontre desfortuno, e quin minacas e frapas}
\l{\cel{3}katastrofi. - II. (\sou{metaf.})\cel{2}Eventajo funesta. - DEFIRS.}
\l{}
\l{\sou{tragikomedio}. Teatrajo di qua la roli ne esas monarki, suve-}
\l{\cel{3}reni, quankam la temo esas tragediala. - DEFIRS.}
\l{}
\l{\sou{trahizar}. (\sou{trans.}) Abandonar (ulu, a qua onu devas esar fi-}
\l{\cel{3}dela), plu o min perfide. - EFS.}
\l{}
\l{\sou{traino}. (\sou{milit-arto}).Korpo di qua la tasko konsistas ye duk-}
\l{\cel{3}tar, transportar, la provizuri, la+asiejo-mashini e la}
\l{\cel{3}municioni. -}
\l{}
\l{\sou{traito}. I. Singla ek la linei di la vizajo. - II. Marko signi-}
\l{\cel{3}fikiva di ulo (karaktero, prudenteso, jenerozeso, e c.)}
\l{\cel{3}- EFI.}
\l{}
\l{\sou{trajektorio}. I. Lineo quan parkuras la baricentro di korpo}
\l{\cel{3}qua lansesis. - II. Kurvo qua krucumas, ye angulo donita,}
\l{\cel{3}kurvi altra, infinite-multa. - DEFIRS.}
\l{}
\l{\sou{trakar.}\cel{2}(\sou{trans.}) I. Persequar (la vildo qua esas en bosko}
\l{\cel{3}o foresto), cirkondante lu, diminutante gradoze la cir-}
\l{\cel{3}klo-diametro. - II. (\sou{metaf.}) Persequar (krimininti, de-}
\l{\cel{3}liktinti) akulante li gradoze a loko sen eskapo-voyo.- F}
\l{}
\l{\sou{trakeo}. (\sou{anatom.)}\cel{1}Tubo ek serio de ringi kartilagoza, qua}
\l{\cel{3}venas de la laringo, decensas alonge la kolo, avan la ezo-}
\l{\cel{3}fago, bifurkas ye lua parto infra, formacante du tubi}
\l{\cel{3}(bronki) di qui singla komunikas kun un ek la du pulmoni,}
\l{\cel{3}ed uzesas por la paso di aero. - DEFIRS.}
\l{}
\l{\sou{trakeito}. (\sou{patol.}) Inflamuro di la trakeo. - DEFIRS.}
\l{}
\l{\sou{trakeotomio}. (\sou{medic.}) Incizo di la trakeo, qua posibligas la}
\l{\cel{3}eniro di la aero extera en la kazo di la afekti qui obs-}
\l{\cel{3}traktas la respiro-tubi (krupo, e c.) - DEFIRS.}
\l{}
\l{\sou{trakito}. (\sou{mineral.}) Roko quan karakterizas la prezenteso}
\l{\cel{3}(plugranda-skale) en olu di feldspato alkalioza, e la}
\l{\cel{3}preska absenteso di quarco o nefelino. - DEFIRS.}
\l{}
\l{\sou{trakomo}. (\sou{patol.})\cel{2}Morbo infektiva di la okul-sistemo, di}
\l{\cel{3}qua la jermo esas ankore nekonocata. - F.}
\l{}
\l{\sou{traktar}.\cel{2}(\sou{trans.}) I. Agar ad (ulu) ye ta o ca maniero. -}
\l{\cel{3}II. (\sou{aludante oratoro, literaturisto, piktisto, e}\cel{1}c.)}
\l{\cel{3}Expozar, developar (tota temo). - EFIRS.}
\l{}
\l{\sou{traktato}. (\sou{pri).}\cel{2}Verko en qua onu traktas temo, explorante}
\l{\cel{3}ye maniero ordinoza lua parti. -}
}
